\section{Introduction}
\label{sec:intro}

A meeting interpreter renders each contribution into the listener's language as it unfolds and makes clear whose
contribution it is: a ``yes'' from the project manager means something different from a ``yes'' from the client. And
people talk at once: over a third of our AMI test windows~\citep{carletta2005ami} hold several speakers.
Simultaneous speech-to-speech translation (SimulS2ST) systems translate one speaker and can keep that speaker's
voice~\citep{seamless2023,labiausse2025hibiki,cheng2025liveinterpret,deng2025simuls2s}. DiariST~\citep{yang2024diarist}
streams speaker-attributed speech-to-text translation of overlapping conversations, without timestamps or speech;
X-Translator~\citep{zhao2026xtranslator} is a speaker-aware speech-to-speech cascade that recognises one stream and routes
each committed span to one speaker, so a merged second voice cannot be recovered. No multi-speaker system reports
latency on overlapping meetings (Appendix~\ref{app:related7}, Table~\ref{tab:related8}).

We present \sys{} (\emph{Meeting Interpretation with Speaker Tracking}), in which one decoder writes, for every speaker
of a window, speaker-attributed, time-stamped translation segments $\spk{k}\,\tim{a}\,\text{translation}\,\tim{b}$,
which a speech generator speaks in the voice of speaker $k$ (Figure~\ref{fig:arch}). The decoder reads a frozen
streaming diarizer, whose slots stay fixed for the whole meeting, as speaker tokens interleaved with its acoustic tokens
(from G-STAR;~\citealp{peng2026gstar}), so two concurrent voices stay two speakers; it writes speaker first, because a
segment's voice must be known before its first word is spoken.

The other half of the problem is \emph{when} to write and to speak. Our first simultaneous decoder learned a wait token
from deadlines, each translated chunk due one fixed lag after the source it covers. It commits only
\num{v10:dev:z:b:pre}\% of its Zh$\rightarrow$En words before the source utterance has ended, and a rule-based speech
layer that plays every committed block in full adds a queue (speech LAAL \num{v10:dev:z:b:slaal}\,s on the development
set). We learn both decisions as latent structure, as two complementary alignments: \emph{write alignment} decides what
text is written and when, \emph{playback alignment} when it is heard. For the decoder, the read at
which each chunk is written (on time or one read later) and the speaker order within a read are latent, and EM
re-weights the training rows by their exact posterior over sets of written chunks. For the speech
generator, each block's rendering and the playback backlog form a linear-chain CRF whose potentials contain the delay at
which each word is heard, under a latency budget imposed by posterior regularisation~\citep{ganchev2010posterior}. More
audio makes the same translation more likely, so EM drifts towards later writes; one constant bias on the wait logit,
set once on development meetings, corrects this at decoding.

Our contributions are (i) the speaker-first, time-stamped joint decoding format with the diarizer read inside the
decoder; (ii) an overlap-aware evaluation with cascade and relabelling controls on the same diarizer
(\S\ref{sec:eval-overlap}); (iii) write alignment by exact EM, with a fixed write bias
(\S\ref{sec:method-latent}); and (iv) playback alignment by exact EM under a latency budget (\S\ref{sec:method-plem}). \sys{} beats the strongest non-oracle cascade on AMI reference windows by
\num{v7:gainoff:en2zh} / \num{v7:gainoff:zh2en} speaker-attributed BLEU (En$\rightarrow$Zh / Zh$\rightarrow$En) and, on
overlapped speech, a speaker-blind model relabelled by the same diarizer by \num{v8:x:ov-O2-rel-en2zh:diff} /
\num{v8:x:ov-O2-rel-zh2en:diff} (\S\ref{sec:res-overlap}). On whole AMI test meetings, against the same joint model
trained on deadlines, the full system raises text tcpBLEU by \num{v10:sig:test:ami:text:m-b:diff} /
\num{v10:sig:test:en:text:m-b:diff} and speech ASR-BLEU by \num{v10:sig:test:ami:speech:m-b:diff} /
\num{v10:sig:test:en:speech:m-b:diff} and changes speech LAAL by \num{v10:test:ami:m-b:slaal} /
\num{v10:test:en:m-b:slaal}\,s (Zh$\rightarrow$En / En$\rightarrow$Zh). Speaking the baseline's own text with playback
alignment yields the speech gains and the latency reduction; write alignment adds the text gain, not speech quality
or speed (\S\ref{sec:res-v10}).
\ifAliwin
On real Mandarin meetings the changes are \num{v10:test:ali:m-b:asr} ASR-BLEU and \num{v10:test:ali:m-b:slaal}\,s.
\else
On real Mandarin meetings the speech gains carry over, but the text score does not improve.
\fi
